\section{A cross-modal probe design for the form--meaning family}\label{sec:D}

This appendix sketches the cross-modal probe that §§\ref{sec:7-2} and \ref{sec:7-4} flag as the open empirical question for the form--meaning family. The aim isn't a finished study but a tractable design whose outcome would settle the cross-domain projectibility question one way or the other.

\subsection{Hypothesis and falsification}

The form--meaning family of §§\ref{sec:3}--\ref{sec:5} comprises three within-domain profiles: theoretical compositionality, rated compound transparency, time-sensitive masked priming. Each licenses within-domain projection. The cross-domain hypothesis is that the three are joint reflexes of a single underlying organization: shared representational structure that exposes form--meaning composability through different windows. The hypothesis is non-trivial because alternative organizations are available: each profile could be domain-internal, with no cross-domain alignment. The probe asks whether the four measurement windows align on the \mention{same items}, which is the only test that distinguishes the cluster reading from the three-separate-profiles reading.

Falsification is built in. If corpus predictors, rated transparency, priming magnitude, and interpretation accuracy fail to align on shared items as joint consequences of representational structure, the family is three within-domain profiles and the cross-domain projectibility question resolves negatively. The diagnostic is a single, item-anchored study; sample-size and statistical-power planning would follow the four-window design rather than any single-window precedent.

\subsection{Item set}

Deadjectival nominalizations on -\mention{ness} are the natural anchor: §\ref{sec:5} establishes a within-domain transparency continuum with rated values from \mention{business} (3.17) to \mention{falseness} (5.75); the -\mention{ness} items have been used in masked-priming work and are corpus-tractable; and parallel Russian -\mention{ost'} items extend the design cross-linguistically without re-doing the architecture. A target set of 60--80 items spanning the rating continuum, with matched controls (frequency-matched simplex nouns, frequency-matched derivations on a different suffix), would support the four windows.

NN compounds (the §\ref{sec:4} target) are an alternative or supplementary anchor. The same within-items logic applies, with Bell and Schäfer's \autocite*{bell-schafer-2016-modelling-transparency} corpus predictors providing the corpus measurement directly. Whether to run the probe on -\mention{ness}, on NN compounds, or on both is a research-design choice; the probe's logical structure is identical across item types.

\subsection{Measurements}

Four within-items measurements are required.

\begin{itemize}
\item
  \term{Corpus predictors}. Frequency, productivity, and semantic-relation expectedness, computed per item from a balanced corpus (BNC, COCA, or COHA depending on the item type). Bell and Schäfer's measurement procedure for compounds extends naturally to nominalizations.
\item
  \term{Rated transparency}. Native-speaker ratings on a 1--7 Likert scale, replicating Heyer and Kornishova's rating procedure with a separate participant group from the priming sub-study.
\item
  \term{Priming magnitude}. Masked-priming reaction times at two SOAs (short, $\approx$30--40 ms; long, $\approx$60--80 ms), with the morphologically related prime, an unrelated simplex control, and an identity prime, replicating Heyer and Kornishova's design within the new item set.
\item
  \term{Interpretation accuracy}. A paraphrase or sentence-completion task probing whether speakers can recover the relation between the derivative and its base in context. The task should be designed so that opaque items predict reduced accuracy and longer response times relative to transparent items.
\end{itemize}

The four windows probe representational structure at distinct stages of processing and explicit access. Their alignment, if found, is a non-trivial empirical pattern.

\subsection{Statistical model}

A joint analysis with item-level random effects and a latent representational-structure factor is the natural statistical move. Each measurement is a noisy reflex; the latent factor's loadings index how strongly each window reflects the underlying structure. The cluster hypothesis predicts non-trivial loadings on all four windows; the three-separate-profiles alternative predicts that the loadings partition by domain, with no shared factor.

The temporal asymmetry §\ref{sec:7-4} flags between priming and interpretation is structural: priming probes early lexical access; interpretation probes explicit recovery; ratings are a metalinguistic judgement supplied at a separate, untimed task. The four windows aren't a temporal chain. The latent-factor model treats them as joint reflexes, with no claim that any one window causes any other.

The design should pre-register a scrambling null alongside the positive test. Permuting items across the four windows and refitting should collapse the shared-factor loadings toward chance; a single factor can otherwise fall out of frequency, familiarity, or method covariance rather than shared representational structure, so only the item-anchored-versus-scrambled contrast distinguishes a genuine cluster from an artefact. The same discipline applies within the priming window, where the SOA-by-transparency interaction is small and should be checked against a label-permutation null. A shared factor that survives item-scrambling, and an interaction that survives label-permutation, are the results that would carry the cross-domain claim.

\subsection{Status}

The probe doesn't exist. The form--meaning family is currently three within-domain profiles plus a research hypothesis. Section \ref{sec:7-4}'s Khalidi-refinement reading is conditional on the probe's outcome; absent the probe, Slater's SPC diagnostic is the relevant standard, and each within-domain profile is treated separately.
